Đồ án nghiên cứu bài toán \textbf{phân bổ tài nguyên vô tuyến (PRB) nhận biết mức độ nguy kịch bệnh nhân} cho nhiều xe cứu thương cùng truyền dữ liệu y tế khẩn cấp (vital signs, cảnh báo DENM/CAM) qua một cell \textbf{5G O-RAN} đô thị. Khi nhiều xe cạnh tranh tài nguyên hữu hạn trong cùng một cell, việc chia đều là không phù hợp về mặt y khoa; tài nguyên cần được ưu tiên theo \textbf{độ nguy kịch} của bệnh nhân trên xe, đồng thời vẫn bảo đảm sàn dịch vụ cho người dùng băng rộng nền (eMBB).

Đóng góp cốt lõi là \textbf{một bài toán tối ưu} được phát biểu dưới dạng \textbf{Quá trình quyết định Markov có ràng buộc (CMDP)}: cực đại hoá tiện ích log-throughput của eMBB \textbf{với năm họ ràng buộc} về trễ trung bình, đuôi trễ (độ tin cậy), sàn thông lượng eMBB, AoI trung bình và đuôi AoI --- mỗi xe chịu ràng buộc theo \textbf{chính mức nguy kịch của nó} (ánh xạ từ thang phân loại ATS 5 mức). Bài toán được giải bằng \textbf{học tăng cường phân cấp hai thang thời gian (HRL)} với tầng Manager (100\,ms, điều khiển ngân sách liên-lát URLLC$\leftrightarrow$eMBB) và tầng Worker (10\,ms, phân chia nội-lát giữa các xe), kết hợp \textbf{đối ngẫu Lagrange} thực thi ràng buộc qua dual ascent. Để bảo đảm so sánh công bằng, cùng một bài toán được giải bởi \textbf{ba solver ngang hàng}: PPO (on-policy), TD3 và SAC (off-policy).

Hệ thống được triển khai trên mô phỏng giải tích (M/G/1 Pollaczek--Khinchine + Monte Carlo + kênh UMa 3GPP TR~38.901) với di động từ vết SUMO/OSM khu vực Bệnh viện Bạch Mai. Tính đúng đắn của \emph{mô hình, công thức và thuật toán} được kiểm chứng bằng \textbf{oracle độc lập} và \textbf{bộ kiểm thử tự động hơn một nghìn ca}; tiêu chí kết luận ``giải đúng'' được định nghĩa chặt chẽ theo \textbf{khả thi sơ cấp có điều kiện theo severity} (kèm quy tắc rule-of-three cho ràng buộc đuôi hiếm), thay vì chỉ dựa trên đường cong phần thưởng.

\medskip\noindent
\textbf{Từ khoá:} O-RAN, network slicing, URLLC, Age of Information, CMDP, Lagrangian, hierarchical reinforcement learning, PPO/TD3/SAC, medical triage.
